% Generated by tools/edn_latex.py from reports/edn.md.
% Do not edit: change reports/edn.md and run `make edn`.
\ednentry{
  id = {EDN-47},
  title = {The serving runtime is \texorpdfstring{\texttt{[project] dependencies}}{[project] dependencies}, and everything else is a PEP 735 dependency group},
  date = {2026-10-05},
  milestone = {M3: Quality Assurance (shapes M5: Model Packaging)},
  topic = {Containerization, Dependency management},
  participants = {@lukas2510},
  decision = {\texttt{pyproject.\allowbreak{}toml} no longer has one flat dependency list.
\texttt{[project] dependencies} holds the serving runtime only, meaning what it takes to load a bundle from \texttt{models/\allowbreak{}<variant>/\allowbreak{}} and price a listing: joblib, lightgbm, loguru, numpy, pandas, pyarrow, python-dotenv, pyyaml, scikit-learn and typer.
Everything else is in \texttt{[dependency-\allowbreak{}groups]}: \texttt{pipeline} (dvc, matplotlib, mlflow, requests, tqdm), \texttt{notebook} (ipython, jupyterlab, notebook), \texttt{docs} (mkdocs), \texttt{test} (coverage, pytest, pytest-cov) and \texttt{dev} (nbstripout, pip, pre-commit, ruff).
\texttt{[tool.\allowbreak{}uv] default-\allowbreak{}groups =\allowbreak{} \textquotedbl{}all\textquotedbl{}} keeps a plain \texttt{uv sync} and \texttt{uv run} installing every group, so no documented command changed.
The API image will install the runtime alone with \texttt{uv sync -\allowbreak{}-\allowbreak{}locked -\allowbreak{}-\allowbreak{}no-\allowbreak{}default-\allowbreak{}groups}.
FastAPI, uvicorn and pydantic join the runtime with the API in M4; Great Expectations, CodeCarbon and \texttt{shap} (EDN-11) go into groups.
\texttt{make test-\allowbreak{}serving}, run by CI's \texttt{Serving runtime} job, which is a required check, builds the runtime set plus the \texttt{test} group from the lock and runs \texttt{tests/\allowbreak{}test\_\allowbreak{}serving.\allowbreak{}py} from it, without \texttt{tests/\allowbreak{}conftest.\allowbreak{}py}.
A serving module that imports a group's package therefore fails the pull request instead of M5's image build.},
  rationale = {\emph{Alternatives considered:}
\begin{itemize}[leftmargin=*, topsep=2pt]
\item \textbf{Option A (chosen): PEP 735 dependency groups, with the runtime in \texttt{[project] dependencies} and \texttt{default-\allowbreak{}groups =\allowbreak{} \textquotedbl{}all\textquotedbl{}}.}
\newline Pros: one lock file and one resolution for everything, so the full set's versions do not move and the image installs a strict subset of exactly what CI tested; uv supports it natively (\texttt{-\allowbreak{}-\allowbreak{}group}, \texttt{-\allowbreak{}-\allowbreak{}only-\allowbreak{}group}, \texttt{-\allowbreak{}-\allowbreak{}no-\allowbreak{}default-\allowbreak{}groups}); with \texttt{default-\allowbreak{}groups =\allowbreak{} \textquotedbl{}all\textquotedbl{}}, a contributor's \texttt{uv sync}, CI's \texttt{uv sync -\allowbreak{}-\allowbreak{}locked} and every \texttt{uv run} behave as before; groups are not published as package metadata, which suits a course project that nobody installs from an index.
\newline Cons: the default is the opposite of what the image needs, so the image and the serving check have to remember \texttt{-\allowbreak{}-\allowbreak{}no-\allowbreak{}default-\allowbreak{}groups}, and a \texttt{uv run} without \texttt{-\allowbreak{}-\allowbreak{}no-\allowbreak{}sync} in that environment would quietly install every group into it.
Two tests in the check close that gap: one fails when the environment holds anything beyond the runtime and \texttt{test} closures read from \texttt{uv.\allowbreak{}lock}, so the check cannot pass vacuously; the other runs the serving path in a fresh interpreter and fails when it imported a module only the \texttt{test} group installs, because the check's environment has pytest and five packages it pulls in, which the image will not.
\item \textbf{Option B: optional extras (\texttt{[project.\allowbreak{}optional-\allowbreak{}dependencies]}).}
\newline Pros: older and more widely known than dependency groups, and every installer has supported them for years.
\newline Cons: \texttt{uv sync} leaves extras out unless it is given \texttt{-\allowbreak{}-\allowbreak{}all-\allowbreak{}extras} or \texttt{-\allowbreak{}-\allowbreak{}extra}, so every documented \texttt{uv sync} (getting-started, CONTRIBUTING, the README, the pipeline docs and every CI job) would either have to change or would silently install less than it does today.
Extras are also published package metadata, an install-time choice for the users of a library, and nobody installs this project as a library.
\item \textbf{Option C: keep one flat list and swap \texttt{mlflow} for \texttt{mlflow-\allowbreak{}skinny}.}
\newline Pros: a one-line change; MLflow is the dependency people usually suspect of bloat.
\newline Cons: measured while reviewing \#42, it does not do what it promises.
\texttt{mlflow-\allowbreak{}skinny} depends on fastapi, uvicorn, starlette and databricks-sdk as well, so none of them leaves.
What it does drop (alembic, docker, flask, flask-cors, graphene, gunicorn, waitress, sqlalchemy, huey, skops, aiohttp) is 39 MB of a 915 MB \texttt{.\allowbreak{}venv}, which is noise next to jupyterlab and the rest of the training stack.
Dropping sqlalchemy and alembic also breaks the two tests that exercise a real MLflow store offline, which are the only place the suite does so without credentials.
And the image would still contain the notebooks, the docs and the tooling.
\item \textbf{Option D: keep one flat list, and give the image its own hand-written requirements file.}
\newline Pros: no change to \texttt{pyproject.\allowbreak{}toml}.
\newline Cons: a second dependency list that nothing ties to \texttt{uv.\allowbreak{}lock}, so the image would either resolve its own versions, and then it is not running what CI tested, or need its versions copied by hand and kept in step forever.
A serving import of a training package would still first fail in the image build.
\item \textbf{Option E: a separate project or uv workspace member for the API.}
\newline Pros: the hardest boundary, with its own \texttt{pyproject.\allowbreak{}toml}.
\newline Cons: the API imports \texttt{recommenditos} itself (\texttt{model.\allowbreak{}py}, \texttt{build\_\allowbreak{}features}), so it would depend on the package anyway and the split would end up inside it again.
A workspace means two members to version and document, which is more structure than a five-person, one-semester project needs.
\end{itemize}
\par
\emph{Why:} The issue framed the lever correctly: what matters is which dependencies reach the image at all, not which MLflow package is installed.
Option C was ruled out by measurement before this work started.
Between A and B the deciding difference is the default.
Extras would have changed every documented \texttt{uv sync}, or quietly installed less with it, while groups with \texttt{default-\allowbreak{}groups =\allowbreak{} \textquotedbl{}all\textquotedbl{}} leave every existing command exactly as it was and only ask the image to opt out.
\par
The runtime set comes from evidence, not from a guess.
An import hook traced a real \texttt{load\_\allowbreak{}model} and \texttt{predict\_\allowbreak{}eur} on every committed bundle, in the full environment and in the runtime one (\texttt{reports/\allowbreak{}analysis/\allowbreak{}serving\_\allowbreak{}imports.\allowbreak{}py}).
\texttt{model.\allowbreak{}py} reaches \texttt{build\_\allowbreak{}features}, \texttt{split\_\allowbreak{}data}, \texttt{pipeline}, \texttt{schema} and \texttt{config}, and those import nine of the ten packages above themselves.
The tenth, pyarrow, is imported by pandas, whose \texttt{read\_\allowbreak{}parquet} needs it but which declares it only as an extra.
The other packages that got imported belong to those ten: their own requirements, such as narwhals, which both lightgbm and scikit-learn require, scipy and threadpoolctl.
The exceptions are optional imports that are skipped when the package is absent: tqdm from \texttt{config.\allowbreak{}py}, psutil from joblib and charset\_normalizer from numpy.
All three appear in the full environment's trace and are missing from the runtime one's, which still serves.
Nothing on the serving path imports MLflow, DVC or anything else from the training stack, so no refactor was needed.
Three small packages stay in the runtime on purpose.
typer is there because every stage module declares its command line when it is imported, and the serving seam imports two of them.
It costs 6.7 MB with rich and Pygments, and splitting \texttt{build\_\allowbreak{}features} to avoid it would break EDN-28's one module per stage and the \texttt{deps} that \texttt{dvc.\allowbreak{}yaml} names.
pyyaml and python-dotenv are there because \texttt{pipeline.\allowbreak{}py} and \texttt{config.\allowbreak{}py} import them, and together they are 3.0 MB.
tqdm went the other way, into \texttt{pipeline}: \texttt{config.\allowbreak{}py} sends loguru through \texttt{tqdm.\allowbreak{}write} with colour forced on whenever tqdm can be imported, which would put escape codes into every container log line.
Without it, the API logs to loguru's default stderr sink, which colours only a terminal.
\texttt{test} is a group of its own rather than part of \texttt{dev}, because the serving check installs it on top of the runtime, and every extra package there would be one a serving import could lean on unnoticed.
matplotlib is in \texttt{pipeline}, not \texttt{notebook}, because a DVC stage draws figures too: issue \#38's \texttt{compare-\allowbreak{}energy} stage writes \texttt{reports/\allowbreak{}figures/\allowbreak{}energy-\allowbreak{}vs-\allowbreak{}error.\allowbreak{}png}, and an install of the pipeline alone has to be able to run it.
\par
Measured on 2026-10-05 by building both environments from \texttt{uv.\allowbreak{}lock} (\texttt{reports/\allowbreak{}analysis/\allowbreak{}runtime\_\allowbreak{}footprint.\allowbreak{}py}): the runtime is \textbf{23 distributions, the project included, and 421 MB of site-packages as installed (519 MB with bytecode)}.
The full environment is \textbf{248 distributions and 815 MB (1,054 MB)}.
So the runtime is 52\,\% of the full environment as installed and 49\,\% with bytecode, and the full one, once compiled, would exceed NFR-04's 1 GB cap with its site-packages alone, before a base image or a model.
The runtime leaves 579 MB as installed, or 481 MB compiled, for the base image, FastAPI, the model and the comparables index.
The largest runtime packages are pyarrow (157 MB), scipy (114 MB), numpy (57 MB), pandas (39 MB), scikit-learn (31 MB) and lightgbm (10 MB).
pyarrow is therefore the first place M5 should look if the image runs tight, but it is not a free saving: \texttt{b0}'s bundle carries a Parquet lookup table, and pandas 3 backs its \texttt{str} dtype with pyarrow whenever it is installed, so dropping it changes how every string column of a request is stored.
The four committed bundles price their first five test rows identically in both environments, so the runtime set serves the same numbers rather than merely importing.
\par
The full set's resolution did not move: \texttt{uv export -\allowbreak{}-\allowbreak{}all-\allowbreak{}groups} lists the same 252 packages at the same versions before and after, the project itself aside, and the only lock change is the project's own entry.
The guarantee is made permanent rather than measured once.
Adding \texttt{import mlflow} to \texttt{model.\allowbreak{}py} made \texttt{make test-\allowbreak{}serving} fail with \texttt{ModuleNotFoundError: No module named \textquotesingle{}mlflow\textquotesingle{}}, while the full suite would still have passed.
Adding \texttt{import packaging.\allowbreak{}version} there failed it as well, naming \texttt{packaging} as a module only the \texttt{test} group installs.
A module-level \texttt{import codecarbon} in \texttt{tests/\allowbreak{}conftest.\allowbreak{}py}, which is what issue \#38's PR does, leaves it green, because the check never loads that file.},
  ai = {Information seeking, Alternative generation, Alternative assessment, Solution generation},
  response = {Accepted},
  assessment = {Lukas chose dependency groups over extras before the implementation started, and the \texttt{mlflow-\allowbreak{}skinny} measurement had already been taken while reviewing \#42.
What AI contributed is the evidence that the split is correct and stays correct.
It traced the serving import graph by running it rather than by reading the imports, which is what separated the three optional imports from real requirements.
It built both environments from the lock, priced the committed bundles in each, checked that the lock's resolution was unchanged, and placed every dependency with a stated reason.
That includes the three that did not obviously belong anywhere: typer kept in the runtime rather than splitting the stage modules, tqdm moved out of it because of the forced-colour log sink, and \texttt{test} separated from \texttt{dev}.
It then proved the CI check catches a training import by injecting one, and added the guard that fails the check when anything beyond the runtime and the test group is installed, so a lost \texttt{-\allowbreak{}-\allowbreak{}no-\allowbreak{}default-\allowbreak{}groups} cannot make the job pass vacuously.
An independent AI review of the pull request then found three ways the check could still be wrong, and reproduced each: the check loaded \texttt{tests/\allowbreak{}conftest.\allowbreak{}py}, so \#38's conftest import of CodeCarbon would have turned it red with no serving module at fault; the test group's own packages could have stood in for a missing runtime dependency; and a change to the \texttt{Makefile} alone did not trigger the job.
Each is fixed and covered by a break-it check; the review also caught that narwhals had been called optional when lightgbm and scikit-learn both require it.},
  aievidence = {Claude Code sessions on 2026-10-05 implementing issue \#61 and reviewing PR \#69: the import trace and the footprint run committed under \texttt{reports/\allowbreak{}analysis/\allowbreak{}}, the before/after \texttt{uv export -\allowbreak{}-\allowbreak{}all-\allowbreak{}groups} comparison, the injected imports that turned \texttt{make test-\allowbreak{}serving} red, and the review's reproduction of the conftest failure by merging \#68's head into this branch.},
  evidence = {\href{https://github.com/mlops-2627q1-mds-upc/MLOPS_Recommenditos/issues/61}{issue \#61}; \href{https://github.com/mlops-2627q1-mds-upc/MLOPS_Recommenditos/pull/69}{PR \#69}; \texttt{pyproject.\allowbreak{}toml}; \href{https://github.com/mlops-2627q1-mds-upc/MLOPS_Recommenditos/blob/main/reports/analysis/runtime_footprint.py}{\texttt{reports/\allowbreak{}analysis/\allowbreak{}runtime\_\allowbreak{}footprint.\allowbreak{}py}} and \href{https://github.com/mlops-2627q1-mds-upc/MLOPS_Recommenditos/blob/main/reports/analysis/runtime_footprint_results.txt}{its output}; \href{https://github.com/mlops-2627q1-mds-upc/MLOPS_Recommenditos/blob/main/reports/analysis/serving_imports.py}{\texttt{reports/\allowbreak{}analysis/\allowbreak{}serving\_\allowbreak{}imports.\allowbreak{}py}} and \href{https://github.com/mlops-2627q1-mds-upc/MLOPS_Recommenditos/blob/main/reports/analysis/serving_imports_results.txt}{its output}; \texttt{tests/\allowbreak{}test\_\allowbreak{}serving.\allowbreak{}py}; \texttt{make test-\allowbreak{}serving} and the \texttt{Serving runtime} job in \texttt{.\allowbreak{}github/\allowbreak{}workflows/\allowbreak{}ci.\allowbreak{}yml}; the Dependencies section of \href{https://github.com/mlops-2627q1-mds-upc/MLOPS_Recommenditos/blob/main/CONTRIBUTING.md}{CONTRIBUTING.md}; NFR-04 in \href{https://github.com/mlops-2627q1-mds-upc/MLOPS_Recommenditos/blob/main/docs/docs/specification.md}{the specification}; \hyperref[EDN-08]{EDN-08}, \hyperref[EDN-11]{EDN-11}, \hyperref[EDN-28]{EDN-28}.},
}
